\section{Problem and motivation}

Much of an organisation's data sits in relational databases, and getting answers out of them normally requires writing SQL. Staff who need a number --- how many students failed more than two subjects, which department has the highest average marks --- usually have to ask someone who can write the query. Text-to-SQL systems remove that step: they translate a plain-English question into a query, run it and return the result.

Recent language models do this well on benchmarks. The difficulty is what happens when they are wrong. A text-to-SQL system answers every question, and a wrong table looks just as authoritative as a correct one; for a user who cannot read SQL, silently wrong answers are the central risk. Most strong systems also depend on closed, hosted models, which means sending the database schema to a third party.

CodeS \cite{li2024codes} addresses the second problem: its weights, code and evaluation setup are public, so it runs on local hardware and its numbers can be checked. It does not address the first. This project starts from the observation that CodeS already computes what is needed to know when it is unsure: it produces four candidate queries with beam search and keeps one. If the candidates agree, the answer is more likely to be right; if they disagree, the question is probably ambiguous or hard. Using that agreement to decide when \emph{not} to answer, and asking the user instead, turns a model that is sometimes silently wrong into one that says when it does not know.

\section{Objectives}

\begin{enumerate}
	\item Reproduce the published Spider result of CodeS within the department's tolerance of $\pm$5--10\%, documenting every deviation from the authors' setup.
	\item Establish which parts of the CodeS pipeline matter, through an ablation study.
	\item Propose, implement and statistically test one enhancement: confidence-aware abstention.
	\item Measure how the system holds up beyond the benchmark, on questions worded as users word them and on real-world databases.
	\item Deliver the system as a working desktop application rather than a notebook.
\end{enumerate}

\section{Contributions and scope}

The project contributes a faithful reproduction of fine-tuned CodeS-1B on Spider \cite{yu2018spider} (77.8\% against 77.9\%) with an audited evaluation harness; an ablation of the CodeS prompt; confidence-aware abstention with a pre-registered hypothesis and significance tests; a benchmark-to-reality evaluation on Spider-Realistic and BIRD; and QueryGenie, a desktop application for macOS and Windows.

CodeS is pre-trained on eight A800 80\,GB GPUs, far beyond a student budget, so the project reproduces the published \emph{inference-time} result with the released fine-tuned checkpoint, which is the configuration behind the 77.9\% figure. Pre-training and fine-tuning were not repeated.
